\begin{solapartat}
\punts{0,2} $A = 0,45-0,30=0,15$ m.

\punts{0,2} $\omega = \sqrt{\dfrac{k}{m}} = \sqrt{\dfrac{300}{2}} = 12,2$ rad/s.

\punts{0,2} $T = \dfrac{2\pi}{\omega} = 0,513$ s

\destaca{Primera opció}

\punts{0,2} Equació del MHS: $x(t) = A\sin(\omega t + \varphi_0)$

\punts{0,2} Inicialment $x_0 = x(t=0) = A$.
\[ A = A\sin(\varphi_0) \Rightarrow \varphi_0 = \mathrm{ArcSin}(1) = \pi/2\ \mathrm{rad}. \]

\punts{0,25} Finalment l'equació del moviment és:
\[ x(t) = 0,15\sin\left(12,2t + \frac{\pi}{2}\right),\ x \text{ en m i } t \text{ en s.} \]

\destaca{Alternativament:}

\punts{0,2} Equació del MHS: $x(t) = A\cos(\omega t + \varphi_0)$

\punts{0,2} Inicialment $x_0 = x(t=0) = A$.
\[ A = A\cos(\varphi_0) \Rightarrow \varphi_0 = \mathrm{ArcCos}(1) = 0. \]

\punts{0,25} Finalment l'equació del moviment és:
\[ x(t) = 0,15\cos(12,2t),\ x \text{ en m i } t \text{ en s.} \]
\end{solapartat}

\begin{solapartat}
\punts{0,25} $E_m = \dfrac{1}{2}kA^2 = 3,38$ J, constant atès que no hi ha dissipació
d'energia

\punts{0,3} $v(t) = \dfrac{dx(t)}{dt} = A\omega\cos(\omega t + \varphi_0) = 1,84\cos(12,2t)$
$v$ en m/s i $t$ en s.
\[ E_c = \frac{1}{2}mv^2 = 3,38\cos^2(12,2t)\ \mathrm{J} \]

\punts{0,3} $U_k = \dfrac{1}{2}kx^2 = 3,38\sin^2(12,2t)\ \mathrm{J}$

\figura[0.55]{sol_fig1.png}

\punts{0,2} Escalat eixos correcte, corbes correctament representades

\punts{0,2} Títols d'eixos amb unitats.
\end{solapartat}
